\section{Immutable Parent Pointers}
\label{sec:immutable-parent-pointers}

TheRecursive Spawning framework establishes ancestry as an immutable, verifiable record. Once a parent pointer is assigned at the moment of spawn, it becomes a permanent, cryptographically sealed fact that no subsequent layer, agent, or process may alter. This immutability is not a courtesy or a policy—it is a structural invariant enforced at theFact Layer. This section defines the formal semantics of parent pointer immutability, the chain traversal algorithm, the enforcement mechanisms, and the spawn authorization protocol that governs pointer assignment.

\subsection{Formal Definition of ParentPointer}

We define the parent pointer relation as a total function over the set of spawned entities.

\begin{definition}[Parent Pointer]
Let $S$ be the set of all spawned entities in the framework. A \textbf{parent pointer} is a binary relation
\[
ParentPointer : S \times S \cup \{\text{null}\}
\]
that maps each entity $e \in S$ to exactly one parent entity $p \in S \cup \{\text{null}\}$. For any $e \in S$:
\[
ParentPointer(e) = 
\begin{cases}
p & \text{if } p \in S \text{ is the sole authorized parent of } e \\
\text{null} & \text{if } e \text{ is the root human (no predecessor)}
\end{cases}
\]
\end{definition}

\begin{invariant}[Uniqueness of Parent]
No entity $e \in S$ may have more than one assigned parent at any point in the framework's existence. Formally:
\[
\forall e \in S, \nexists p_1, p_2 \in S \cup \{\text{null}\} : p_1 \neq p_2 \land ParentPointer(e) = p_1 \land ParentPointer(e) = p_2
\]
This uniqueness is not enforced dynamically—it is immutably fixed at the moment of spawn.
\end{invariant}

\begin{invariant}[Immutability of Assigned Pointer]
Once $ParentPointer(e)$ is assigned a value $p \neq \text{null}$ at the time of spawn, the relation $ParentPointer(e) = p$ is permanently sealed. Formally:
\[
\forall e \in S, \forall t > t_{\text{spawn}}(e) : ParentPointer(e) = p \quad \text{is constant and unchangeable}
\]
No agent, Purpose Layer directive, or external authorization may modify $ParentPointer(e)$ after $t_{\text{spawn}}(e)$.
\end{invariant}

The implication is that ancestry is not a mutable property that can be edited, reassigned, or corrected after the fact. If a spawn event is recorded with a parent pointer to $p$, that pointer remains pointing to $p$ for the entire lifetime of the entity and the framework.

\subsection{The Chain Relation}

Given the immutability of individual parent pointers, the full ancestry trace of any entity is defined recursively as the \textbf{Chain}.

\begin{definition}[Ancestry Chain]
For any entity $a \in S$, the ancestry chain $\Chain(a)$ is defined recursively as:
\[
\Chain(a) \equiv ParentPointer(parent(a), a) \cup \Chain(parent(a))
\]
where $parent(a)$ is the unique entity such that $ParentPointer(a) = parent(a)$.

The base case is:
\[
\Chain(r) = \emptyset \quad \text{for the root human } r \in S \text{ where } ParentPointer(r) = \text{null}
\]
\end{definition}

In words, the chain of $a$ consists of the edge connecting $a$ to its parent, plus the entire chain of its parent, all the way back to the root. Since each $ParentPointer$ is immutable, $\Chain(a)$ is a monotonically growing set—once an edge is in the chain it can never be removed or altered.

\begin{figure}[H]
\centering
\begin{tikzpicture}[
    node distance=2.5cm,
    every node/.style={circle, draw, minimum size=1.2cm, font=\footnotesize}
]
    \node (r) {$r$\\$\emptyset$};
    \node (p) [right of=r] {$p$\\$\Chain(p)$};
    \node (a) [right of=p] {$a$\\$\Chain(a)$};
    \node (null) [below of=r, node distance=1.5cm] {null parent};

    \draw[-Stealth] (p) -- (r) node[midway, above] {$ParentPointer(p) = r$};
    \draw[-Stealth] (a) -- (p) node[midway, above] {$ParentPointer(a) = p$};

    \node [draw=none, font=\footnotesize, align=left, anchor=north west] at (current bounding box.south west) {
        $\Chain(p) = \{ParentPointer(r, p)\}$ \\
        $\Chain(a) = \{ParentPointer(p, a)\} \cup \Chain(p) = \{ParentPointer(p, a), ParentPointer(r, p)\}$
    };
\end{tikzpicture}
\caption{Recursive chain construction: $\Chain(a)$ extends $\Chain(parent(a))$ by one edge.}
\label{fig:chain-recursion}
\end{figure}

The chain relation supports efficient traversal through the following algorithm.

\begin{algorithm}[Chain Traversal]
Given an entity $a$, the chain traversal algorithm enumerates all ancestors in order from $a$ to the root:

\begin{enumerate}
    \item Initialize $\cursor \leftarrow a$, $\chain \leftarrow \emptyset$.
    \item While $ParentPointer(\cursor) \neq \text{null}$:
    \begin{enumerate}
        \item Record the edge: $\chain \leftarrow \chain \cup \{(\cursor, ParentPointer(\cursor))\}$.
        \item Advance: $\cursor \leftarrow ParentPointer(\cursor)$.
    \end{enumerate}
    \item Return $\chain$, which is exactly $\Chain(a)$.
\end{enumerate}

The algorithm terminates because $S$ is finite and each step moves to a strictly higher ancestor; the recursion depth is bounded by $|S|$.
\end{algorithm}

\begin{note}
The chain traversal algorithm is read-only. It never modifies any parent pointer—it only reads the immutable relations to construct the ancestry record. This is a critical property: traversal does not compromise immutability.
\end{note}

\subsection{Immutability Enforcement at the Fact Layer}

The Fact Layer is the enforcement substrate for parent pointer immutability. It does not merely store the pointer—it guarantees that no modification can occur through any pathway.

\begin{invariant}[Fact Layer Write Protection]
At the moment of spawn, theFact Layer performs the following actions atomically:

\begin{enumerate}
    \item \textbf{Assignment.} $ParentPointer(c) \leftarrow p$ is recorded in the factual store with a timestamp $t_{\text{spawn}}(c)$.
    \item \textbf{Sealing.} The record is cryptographically hashed and sealed. The hash $h = \text{hash}(c, p, t_{\text{spawn}}(c))$ is stored alongside the relation, forming an tamper-evident log entry.
    \item \textbf{Lock.} A write-protection flag is set on the $ParentPointer$ field of entity $c$. From $t > t_{\text{spawn}}(c)$ onward, all write operations to $ParentPointer(c)$ are rejected at the Fact Layer.
\end{enumerate}

Formally, theFact Layer enforces:
\[
\forall e \in S, \forall t > t_{\text{spawn}}(e) : \text{write\_attempt}(ParentPointer(e), t) \rightarrow \text{REJECTED}
\]
\end{invariant}

\begin{figure}[H]
\centering
\begin{tikzpicture}[
    node distance=3cm,
    block/.style={rectangle, draw, minimum height=1.2cm, text width=2.8cm, align=center, font=\footnotesize},
    arrow/.style={-Stealth, thick}
]
    \node [block] (spawn) {Spawn Event \\ $ParentPointer(c) \leftarrow p$};
    \node [block, right of=spawn] (factlayer) {Fact Layer \\ Atomic Write + Seal};
    \node [block, right of=factlayer] (locked) {Immutable Record \\ Write-Protected};

    \draw[arrow] (spawn) -- (factlayer);
    \draw[arrow] (factlayer) -- (locked);

    \node [block, below of=factlayer, node distance=2cm] (reject) {Modification Attempt \\ $\rightarrow$ REJECTED};

    \draw[dashed, arrow] (reject) -- (factlayer);
\end{tikzpicture}
\caption{Fact Layer enforces immutability at spawn: pointer is assigned, sealed, and locked in a single atomic operation.}
\label{fig:fact-layer-immutability}
\end{figure}

The Fact Layer treats any attempted modification of an assigned parent pointer as a category-one integrity violation. The framework does not issue a warning or a soft rejection—it treats the modification attempt as a structural fault. This is consistent with the principle that ancestry is not supplementary metadata but a core structural property of every entity in the system.

\subsection{Spawn Authorization and the Purpose Layer}

While theFact Layer enforces immutability after assignment, the \textbf{Purpose Layer} governs whether a spawn event may occur in the first place—and therefore whether a parent pointer may be created. The Purpose Layer is the only authoritative source that may approve the creation of a new $ParentPointer$ relation.

\begin{protocol}[Spawn Authorization for ParentPointer Assignment]
When a spawn event is proposed with candidate parent $p$ and candidate child $c$, the Purpose Layer performs the following checks before authorizing the creation of $ParentPointer(c) = p$:

\begin{enumerate}
    \item \textbf{Entity Validity.} Verify that both $p$ and $c$ are registered entities in theFact Layer with confirmed identity.
    \item \textbf{Authorization Scope.} Confirm that the agent or process requesting the spawn holds a valid spawn authorization grant from the Purpose Layer.
    \item \textbf{Prohibition Check.} Verify that no existing $ParentPointer(c)$ record is already present. If $c$ already has a parent assigned, the spawn is denied—the pointer, once created, is immutable and cannot be overwritten.
    \item \textbf{Relation Legality.} Confirm that the proposed parent-child relation does not violate the Purpose Layer's spawn policy (e.g., constraints on lineage depth, authorization scope boundaries).
\end{enumerate}

If all checks pass, the Purpose Layer issues a spawn authorization token $T_{\text{spawn}}$ and forwards the validated $(p, c)$ pair to the Fact Layer for pointer creation.
\end{protocol}

\begin{note}
The Purpose Layer authorizes the creation of the pointer—it does not and cannot modify an already-created pointer. Authorization is a one-time gate at spawn; after the gate closes, theFact Layer's immutability guarantees take over.
\end{note}

The separation of concerns between the two layers is critical:

\begin{itemize}
    \item \textbf{Purpose Layer}: Governs whether a parent pointer \textit{may be created}. It is the gatekeeper of spawn authorization.
    \item \textbf{Fact Layer}: Governs whether a parent pointer \textit{may be modified after creation}. It is the enforcer of immutability.
\end{itemize}

Together, these two layers guarantee that the ancestry record is both \textit{authoritatively created} and \textit{permanently preserved}.

\subsection{Retrieval and Verification}

Given an entity $a$, any participant in the framework—whether another agent, an audit process, or an external verifier—can retrieve and verify the full ancestry chain without ambiguity.

\begin{verification}[Chain Verification]
To verify that a chain $\Chain(a)$ is structurally sound:
\begin{enumerate}
    \item For each edge $(child, parent) \in \Chain(a)$, confirm that $ParentPointer(child) = parent$ is sealed in theFact Layer.
    \item Confirm that the final element of the chain points to the unique root entity $r$ where $ParentPointer(r) = \text{null}$.
    \item Confirm that no entity in the chain (including $a$) appears more than once, which would indicate a cycle—an impossibility given the one-parent-per-entity constraint.
\end{enumerate}
Because each $ParentPointer$ is immutable and sealed, verification requires only a read operation against theFact Layer. There is no need for forensic reconstruction or conflict resolution—the record is definitive.
\end{verification}

\subsection{Properties of the Immutable Chain}

The combination of immutable parent pointers and recursive chain construction yields a set of formally guaranteed properties.

\begin{property}[No Cycles]
The chain of any entity is acyclic. Formally:
\[
\forall a \in S, \nexists \text{ sequence } a_0, a_1, \ldots, a_k \text{ s.t. } a_0 = a_k \land \forall i : ParentPointer(a_{i+1}) = a_i
\]
This holds because each entity has at most one parent, and the chain traversal strictly ascends toward null. A cycle would require an entity to be its own ancestor, which is impossible under the one-parent-per-entity constraint.
\end{property}

\begin{property}[Deterministic Ancestry]
For any entity $a$, $\Chain(a)$ is a deterministic, uniquely defined set. Given the same entity $a$, any traversal algorithm will produce exactly the same ancestry record. There is no ambiguity, no divergence, and no competing interpretations of ancestry.
\end{property}

\begin{property}[Monotonic Growth of Chain Information]
As the framework evolves and new entities spawn, the union of all chains across all entities grows monotonically—it never contracts or rewrites. This means the entire ancestry graph is an append-only, tamper-evident ledger of the complete spawning history of the system.
\end{property}

These properties ensure that the ancestry record produced by theRecursive Spawning framework is trustworthy without adjudication. There is no need for a disputed claim resolution mechanism, because theFact Layer's immutability guarantees eliminate the possibility of disputed ancestry in the first place.

\begin{table}[H]
\centering
\caption{Summary of Immutable Parent Pointer Properties}
\label{tab:parent-pointer-properties}
\begin{tabular}{@{}L{4cm}L{6cm}L{4cm}@{}}
\toprule
\textbf{Property} & \textbf{Formal Statement} & \textbf{Enforcement} \\
\midrule
Uniqueness & $\forall e \in S, \exists !$ $p : ParentPointer(e) = p$ & Purpose Layer at spawn \\
\midrule
Immutability & $\forall e \in S, \forall t > t_{\text{spawn}}(e) : ParentPointer(e) = p$ is constant & Fact Layer write lock \\
\midrule
Chain Recursion & $\Chain(a) = ParentPointer(parent(a), a) \cup \Chain(parent(a))$ & Traversal algorithm \\
\midrule
Base Case & $\Chain(r) = \emptyset$ for $r$ where $ParentPointer(r) = \text{null}$ & Definition \\
\midrule
Acyclicity & $\nexists$ cycle in $\Chain(a)$ for any $a \in S$ & One-parent constraint \\
\midrule
Tamper-Evidence & Each pointer sealed with cryptographic hash at $t_{\text{spawn}}$ & Fact Layer sealing \\
\bottomrule
\end{tabular}
\end{table}

The immutable parent pointer system, therefore, provides the foundational structural guarantee of theRecursive Spawning framework: every spawned entity has a permanent, verifiable, unalterable record of its ancestry, from its first authorized spawn all the way back to the root human.